% Part: first-order-logic
% Chapter: models-theories
% Section: set-theory

\documentclass[../../../include/open-logic-section]{subfiles}

\begin{document}

\olfileid{fol}{mat}{set}

\olsection{संचांची उपपत्ती}

जवळजवळ संपूर्ण गणित संचांच्या उपपत्तीमध्ये विकसित करता येते. या
उपपत्तीमध्ये गणित विकसित करताना अनेक गोष्टी कराव्या लागतात. प्रथम,
संबंध~$\in$ साठी स्वयंसिद्धकांचा संच लागतो. वेगवेगळ्या अनेक
स्वयंसिद्धक-प्रणाली विकसित झाल्या असून, त्यांतील~$\in$ चे गुणधर्म कधीकधी
परस्परविरोधी असतात. निवडीच्या स्वयंसिद्धकासह झर्मेलो--फ्रेंकेल संच
उपपत्ती म्हणून ओळखली जाणारी स्वयंसिद्धक-प्रणाली~$\Log{ZFC}$ विशेष उठून
दिसते: ती आतापर्यंत सर्वाधिक वापरली आणि अभ्यासली गेलेली प्रणाली आहे, कारण
गणितज्ञांना ज्या जवळजवळ सर्व गोष्टी सिद्ध करता याव्यात अशी अपेक्षा असते,
त्या सिद्ध करण्यास तिची स्वयंसिद्धके पुरेशी ठरतात. पण हे स्थापित करण्यापूर्वी
गणितज्ञांना व्यक्त करायच्या सर्व गोष्टी आपण \emph{व्यक्त} तरी कशा करू शकतो,
हे प्रथम स्पष्ट करणे आवश्यक आहे. सुरुवातीलाच, भाषेत कोणतीही !!{constant}s
किंवा !!{function}s नसतात; त्यामुळे विशिष्ट संचांविषयी (उदाहरणार्थ,
$\emptyset$ किंवा~$\Nat$), संचांवरील क्रियांविषयी (उदाहरणार्थ,
$X \cup Y$ आणि~$\Pow{X}$), इतकेच नव्हे तर संबंध आणि फलने यांसारख्या
संचांखेरीज इतर गोष्टींचा समावेश असलेल्या रचनांविषयी आपण बोलू शकतो हेही
पहिल्या नजरेत अस्पष्ट वाटते.

सुरुवातीला, ``चा घटक आहे'' हाच काही आपल्याला महत्त्वाचा असलेला एकमेव संबंध
नाही; ``चा उपसंच आहे'' हा जवळजवळ तितकाच महत्त्वाचा वाटतो. पण ``चा घटक
आहे'' याच्या आधारे ``चा उपसंच आहे'' हे आपण \emph{परिभाषित} करू शकतो. त्यासाठी
संच उपपत्तीच्या भाषेत असे !!a{formula}~$!A(x, y)$ शोधावे लागते की संचांची
जोडी~$\tuple{X, Y}$ त्याची पूर्ति तेव्हा आणि केवळ तेव्हाच करते, जेव्हा
$X \subseteq Y$. पण~$X$ हा~$Y$ चा उपसंच तेव्हा आणि केवळ तेव्हाच असतो,
जेव्हा~$X$ चे सर्व !!{element}s हे~$Y$ चेही !!{element}s असतात. म्हणून
पुढील सूत्राने~$\subseteq$ ची व्याख्या करता येते:
\[
\lforall[z][(z \in x \lif z \in y)]
\]
आता, एखाद्या सूत्रात~$\subseteq$ संबंध वापरायचा असेल तेव्हा त्याऐवजी हे
सूत्र वापरता येईल ($x$ आणि~$y$ यांच्या जागी योग्य पदे घालून आणि आवश्यक
असल्यास बद्ध चर~$z$ चे नाव बदलून). उदाहरणार्थ, संचांची विस्तारात्मकता
म्हणजे कोणतेही संच~$x$ आणि~$y$ हे परस्परांचे उपसंच असतील, तर~$x$ आणि~$y$
एकच संच असले पाहिजेत. हे
$\lforall[x][\lforall[y][((x \subseteq y \land y
    \subseteq x) \lif x = y)]]$ ने व्यक्त करता येते; किंवा वरच्या व्याख्येने
$\subseteq$ बदलल्यास पुढील सूत्राने:
\[
\lforall[x][\lforall[y][((\lforall[z][(z \in x \lif z \in y)] \land
    \lforall[z][(z \in y \lif z \in x)]) \lif x = y)]].
\]
हे खरे तर~$\Log{ZFC}$ च्या स्वयंसिद्धकांपैकी एक, ``विस्तारात्मकतेचे
स्वयंसिद्धक,'' आहे.

$\emptyset$ साठी कोणतेही !!{constant} नाही, पण ``$x$ रिक्त आहे'' हे
$\lnot \lexists[y][y \in x]$ ने व्यक्त करता येते. मग ``$\emptyset$
अस्तित्वात आहे'' हे !!{sentence}~$\lexists[x][\lnot \lexists[y][y \in
    x]]$ बनते. हे~$\Log{ZFC}$ चे आणखी एक स्वयंसिद्धक आहे. (विस्तारात्मकतेच्या
स्वयंसिद्धकानुसार फक्त एकच रिक्त संच असतो, हे लक्षात घ्या.) संच उपपत्तीच्या
भाषेत~$\emptyset$ विषयी बोलायचे असेल तेव्हा आपण ते ``रिक्त असलेला एक संच
अस्तित्वात आहे आणि \dots'' असे लिहू. उदाहरणार्थ, $\emptyset$ हा प्रत्येक संचाचा
उपसंच आहे ही गोष्ट पुढीलप्रमाणे व्यक्त करता येईल:
\[
\lexists[x][(\lnot\lexists[y][y \in x] \land \lforall[z][x \subseteq
    z])]
\]
येथे अर्थातच $x \subseteq z$ हेसुद्धा त्याच्या व्याख्येने बदलावे लागेल.

$X \cup Y$ आणि~$\Pow{X}$ यांसारख्या संचांवरील क्रियांविषयी बोलण्यासाठी
अशीच युक्ती वापरावी लागते. संच उपपत्तीच्या भाषेत फलनचिन्हे नसतात, पण
$X \cup Y = Z$ आणि~$\Pow{X} = Y$ हे फलनीय संबंध पुढील सूत्रांनी व्यक्त
करता येतात:
\begin{align*}
& \lforall[u][((u \in x \lor u \in y) \liff u \in z)]\\
& \lforall[u][(u \subseteq x \liff u \in y )]
\end{align*}
कारण $X \cup Y$ चे !!{element}s म्हणजे नेमके~$X$ चे !!{element}s किंवा
~$Y$ चे !!{element}s असलेले संच, आणि~$\Pow{X}$ चे !!{element}s म्हणजे नेमके~$X$
चे उपसंच होत. पण त्यामुळे आपण $x \cup y$ किंवा~$\Pow{x}$ ही पदे असल्यासारखी
वापरू शकत नाही; $X \cup Y = Z$ आणि~$\Pow{X} = Y$ हे संबंध परिभाषित करणारी
संपूर्ण !!{formula}s एवढीच वापरता येतात. प्रत्यक्षात, या संबंधांची कधी तरी
पूर्ति होते, म्हणजे संयोग आणि घातसंच नेहमी अस्तित्वात असतात, हे आपल्याला
अजून माहीत नाही. उदाहरणार्थ, !!{sentence}
$\lforall[x][\lexists[y][\Pow{x} = y]]$ हे~$\Log{ZFC}$ चे आणखी एक
स्वयंसिद्धक (घातसंचाचे स्वयंसिद्धक) आहे.

आता क्रमित जोड्या किंवा फलने यांविषयी कसे बोलायचे? येथे क्रमित जोड्या आणि
फलने हे विशिष्ट प्रकारचे संच म्हणून कसे मानता येतात, हे स्पष्ट करावे लागते.
क्रमित जोडी~$\tuple{x, y}$ ची व्याख्या संच~$\{\{x\}, \{x, y\}\}$ म्हणून
करणे हा एक मार्ग आहे. पण आधीप्रमाणेच या संचाला नाव देणारे !!a{function}
आपण समाविष्ट करू शकत नाही; फक्त $\tuple{x, y} = z$, म्हणजे
$\{\{x\}, \{x, y\}\} = z$, हा संबंध परिभाषित करता येतो:
\[
\lforall[u][(u \in z \liff (\lforall[v][(v \in u \liff v = x)] \lor
  \lforall[v][(v \in u \liff (v = x \lor v = y))]))]
\]
याचा अर्थ~$z$ चे !!{element}s~$u$ म्हणजे नेमके असे संच की ज्यांचा एकमेव
!!{element} हा~$x$ असतो किंवा ज्यांचे एकमेव !!{element}s हे~$x$ आणि~$y$
असतात (दुसऱ्या शब्दांत, जे संच~$\{x\}$ शी किंवा~$\{x, y\}$ शी एकरूप
असतात). हे मिळाल्यानंतर आणखी गोष्टीही सांगता येतात; उदाहरणार्थ,
$X \times Y = Z$:
\[
\lforall[z][(z \in Z \liff \lexists[x][\lexists[y][(x \in
        X \land y \in Y \land \tuple{x, y} = z)]])]
\]

फलन~$f \colon X \to Y$ याचा संबंध~$f(x) = y$, म्हणजे जोड्यांचा संच
$\Setabs{\tuple{x,y}}{f(x) = y}$, म्हणून विचार करता येतो. मग संच~$f$ हे
$X$ पासून~$Y$ पर्यंतचे फलन आहे असे आपण तेव्हा म्हणू शकतो, जेव्हा (a) तो
$\subseteq X \times Y$ असा संबंध आहे, (b) तो सर्वत्र परिभाषित आहे,
म्हणजे प्रत्येक $x \in X$ साठी असा काही $y \in Y$ आहे की
$\tuple{x, y} \in f$, आणि (c) तो एकमूल्यी आहे, म्हणजे
$\tuple{x, y}, \tuple{x, y'} \in f$ असताना $y = y'$ (कारण फलनांची मूल्ये
एकमेव असली पाहिजेत). म्हणून ``$f$ हे~$X$ पासून~$Y$ पर्यंतचे फलन आहे'' हे
पुढीलप्रमाणे लिहिता येते:
\begin{align*}
 \lforall[u][(u \in f \lif {}] & \lexists[x][\lexists[y][(x \in X \land y \in
      Y \land \tuple{x, y} = u)]]) \land {}\\
 \lforall[x][(x \in X \lif {}] &
      (\lexists[y][(y \in Y \land \mathrm{maps}(f, x, y))] \land {}\\
& (\lforall[y][\lforall[y'][((\mathrm{maps}(f, x, y) \land
    \mathrm{maps}(f, x, y')) \lif y = y')]]))
\end{align*}
येथे $\mathrm{maps}(f, x, y)$ हे $\lexists[v][(v \in f \land
  \tuple{x, y} = v)]$ चे संक्षिप्त रूप आहे (हे !!{formula} ``$f(x) = y$''
व्यक्त करते).

आता, उदाहरणार्थ, $f\colon X \to Y$ हे !!{injective} आहे हे व्यक्त करणेही
कठीण नाही:
\begin{multline*}
 f \colon X \to Y \land \lforall[x][\lforall[x'][((x \in X \land x' \in
     X \land {}]] \\
 \lexists[y][(\mathrm{maps}(f, x, y) \land \mathrm{maps}(f,
        x', y))]) \lif x = x')
\end{multline*}
फलन~$f\colon X \to Y$ हे !!{injective} तेव्हा आणि केवळ तेव्हाच असते, जेव्हा
$f$ हे $x, x' \in X$ या दोन्हींना एकाच~$y$ शी संगत करत असल्यास $x = x'$.
हे सूत्र~$\mathrm{inj}(f, X, Y)$ असे संक्षिप्तपणे लिहिल्यास, आता आपण संच
उपपत्तीच्या भाषेत कँटरच्या प्रमेयासारखी साधी नसलेली गोष्टही मांडू शकतो:
$\Pow{X}$ पासून~$X$ मध्ये जाणारे कोणतेही !!{injective} फलन नाही:
\[
\lforall[X][\lforall[Y][(\Pow{X} = Y \lif
    \lnot\lexists[f][\mathrm{inj}(f, Y, X)])]]
\]

प्रत्येक परिभाषक गुणधर्मासाठी संचाचे अस्तित्व हमीने देणारे आणखी एक
स्वयंसिद्धक संच उपपत्तीला आवश्यक आहे असे वाटू शकते. $!A(x)$ हे मुक्त
चर~$x$ असलेले संच उपपत्तीचे सूत्र असेल, तर पुढील !!{sentence} विचारात घेता
येते:
\[
\lexists[y][\lforall[x][(x \in y \liff !A(x))]].
\]
हे !!{sentence} सांगते की असा संच~$y$ आहे ज्याचे !!{element}s म्हणजे नेमके
ते सर्व~$x$ होत जे~$!A(x)$ ची पूर्ति करतात. या रूपबंधाला ``गुणधर्माधारित
संचरचनेचे तत्त्व'' म्हणतात. ते फार उपयुक्त वाटते; पण दुर्दैवाने ते विसंगत
आहे. $!A(x) \ident \lnot x \in x$ घ्या; मग गुणधर्माधारित संचरचनेचे तत्त्व
पुढील विधान सांगते:
\[
\lexists[y][\lforall[x][(x \in y \liff x \notin x)]],
\]
म्हणजे स्वतःचे !!{element}s नसलेल्या सर्व संचांच्या संचाचे अस्तित्व ते
सांगते. असा संच अस्तित्वात असू शकत नाही---ही रसेलची विरोधापत्ती आहे.
प्रत्यक्षात, $\Log{ZFC}$ मध्ये या तत्त्वाची मर्यादित---आणि सुसंगत---आवृत्ती,
म्हणजे विभक्तीकरण तत्त्व, समाविष्ट आहे:
\[
\lforall[z][\lexists[y][\lforall[x][(x \in y \liff (x \in z \land
      !A(x))]]].
\]

\begin{prob}
पुढील निष्पन्नता दाखवणारी !!a{derivation} देऊन गुणधर्माधारित संचरचनेचे
तत्त्व विसंगत आहे, हे दाखवा:
\[
\lexists[y][\lforall[x][(x \in y \liff x \notin x)]] \Proves \lfalse.
\]
प्रथम $(A \lif \lnot A) \land (\lnot A \lif A)
\Proves \lfalse$ दाखवल्यास मदत होईल.
\end{prob}
\end{document}
